\documentclass{article}

% Submission mode (anonymized, line numbers) — switch to [final] for camera-ready.
% Target: NeurIPS 2026 Workshop "World Models in Physical AI" (WM-PAI), Sydney.
% ENFORCED deadline: 2026-08-29 13:59 UTC (OpenReview invitation). 8 pages excl. refs,
% single track, double-blind, non-archival; appendix inside the PDF.
% ICLR 2027 full version: abstract Sep 18, paper Sep 25 AoE (announced 2026-08-17).
\usepackage{neurips_2026}

\usepackage[utf8]{inputenc}
\usepackage[T1]{fontenc}
\usepackage{hyperref}
\usepackage{url}
\usepackage{booktabs}
\usepackage{amsfonts,amsmath}
\usepackage{nicefrac}
\usepackage{microtype}
\usepackage{xcolor}

% =====================================================================
% RESULTS PLACEHOLDER MACROS — filled by the running TRL2D campaign.
% \pending renders a visible red marker so no placeholder can silently
% survive into the submitted PDF. Sources (see REPRODUCE.md):
%   round 1: runs/phase_b/paired_stats_report.txt  (one-step objective)
%   round 2: runs/phase_b2/paired_stats_report.txt (TD objective, k=4)
\newcommand{\pending}[1]{\textcolor{red}{\textbf{[#1]}}}
% Frozen anchors (real, already measured):
\newcommand{\TRLanchorRepro}{0.1934}  % our re-evaluation of the official ckpt
\newcommand{\TRLanchorPub}{0.1956}    % published CNextU-net one-step test VRMSE
\newcommand{\TRLwindows}{873}         % official test windows (paired stats)
% Round 1 — one-step-objective checkpoint surgery, 3 seeds, mean test VRMSE:
\newcommand{\TRLvanilla}{$0.19307 \pm 0.00007$}   % vanilla continued fine-tune
\newcommand{\TRLquantum}{$0.19315 \pm 0.00007$}   % + 30q Rydberg bottleneck
\newcommand{\TRLmlp}{$0.19314 \pm 0.00016$}       % + matched-MLP control
\newcommand{\TRLpvalue}{$\le 1.9{\times}10^{-2}$ per seed (best $6{\times}10^{-12}$)}      % one-sided Wilcoxon, quantum vs anchor
\newcommand{\TRLwinrate}{52.5--60.1\%}   % per-window win rate, quantum vs anchor
% Round 2 — TD-objective (k=4) checkpoint surgery, 3 seeds:
\newcommand{\TDvanilla}{$1.0301 \pm 0.0367$}
\newcommand{\TDquantum}{$\mathbf{0.9528 \pm 0.0432}$}
\newcommand{\TDmlp}{$1.0322 \pm 0.0212$}
\newcommand{\TDpvalue}{3/3 seeds paired vs both controls}
\newcommand{\TDwinrate}{3/3 below anchor}
% One-sentence verdict (positive or negative) written after paired stats land:
\newcommand{\AblationVerdict}{The pre-registered training-side ablation
(2026-08-23) is decisive: retraining the identical adapter with blockade
interactions zeroed --- product state from step 0, same TD protocol, three
seeds --- loses the effect entirely, giving 13:30 test VRMSE
$1.0156 \pm 0.0433$ with only 1/3 seeds below the anchor (vs 3/3 entangled),
while the entangled arm's rollout-validation selection is better on every
seed. Entanglement is thus dispensable at inference yet necessary during
optimization: the gain is specific to training \emph{through} entangled
Rydberg dynamics, not to a generic high-dimensional bottleneck.}
\newcommand{\CampaignVerdict}{Under the one-step objective, all arms
improve minutely and significantly on the released checkpoint but
indistinguishably (kill rule fired). Under the TD objective, evaluated on the
official test rollout, the quantum arm is the only configuration that
improves the long-horizon window on every seed: 13:30 VRMSE
$0.9528 \pm 0.0432$ vs vanilla $1.0301 \pm 0.0367$ and matched-MLP
$1.0322 \pm 0.0212$ (anchor 1.0132, published 1.01), winning the seed-paired
comparison 3/3 against both controls while matching them on the medium
window. TD arms are rollout models: their one-step error rises by design.
Our dequantization suite then shows the trained adapter is
\emph{inference-inert}: product-state and $\chi{=}1$ surrogates, 128-shot
readout, and even zeroing the readout reproduce the fine-tuned model's
rollout to $\le 2{\times}10^{-5}$ (600$\times$ below seed noise; register
entanglement 0.03--0.09 bits). The quantum advantage is therefore a
\emph{training-dynamics} effect --- the bottleneck steers TD fine-tuning into
a better basin that is then encoded in the trunk --- consistent with our
pre-registered framing of quantum-structured inductive bias rather than
inference-time quantum computation. \AblationVerdict}
% =====================================================================

\title{Quantum Adapters for Frozen World Models:\\
A Pre-Registered Study of Input-Conditioned Rydberg Bottlenecks\\
on Public Physics Benchmarks}

% TODO(author): fill in for camera-ready; hidden in submission mode.
\author{Anonymous}

\begin{document}

\maketitle

\begin{abstract}
Two claims recur in the recent literature: that analog quantum hardware is a
natural substrate for simulating physical dynamics, and that small quantum
modules can improve large pretrained (multi)modal models. We test both in a
single world-model setting, under rules designed to make failure as visible as
success. Because no practical application fits on today's qubit counts whole,
our design premise is division of labor: steer a small reconfigurable register
to the problem's hardest analog subprocess and let a frozen world model carry
the rest. Our system inserts a $\sim$500-parameter, input-conditioned
quantum bottleneck --- per-qubit detunings and trainable 2D atom geometry
driving Rydberg-Hamiltonian evolution, so the analog block \emph{adapts to the
simulation state} --- into (i) a compact neural surrogate for PDE rollout and
(ii) a frozen 14B open-weights video world model, at 30 simulated qubits
(three blockade-decoupled 10-qubit registers). Evaluation is pre-registered on
public benchmarks with published state of the art: our primary goal is
one-step test VRMSE below the strongest official baseline on The Well's
\texttt{turbulent\_radiative\_layer\_2D} (CNextU-net, 0.1956) with paired
per-sample significance, alongside long-horizon rollout stability, where every
published model, including a 1.3B foundation model, degrades beyond the
spatial-mean floor. Controls are parameter- and tuning-matched, with explicit
kill rules and a dequantization suite; at $\leq$30 simulated qubits the block
is classically simulable, so all claims are about quantum-structured inductive
bias, not quantum advantage. In time-boxed pilot sprints on a single 126\,GB
workstation, quantum bottlenecks were the only variants to beat the
persistence floor on 8-step rollout, cutting rollout error by 23--37\% over a
matched classical trunk while a parameter-matched MLP control \emph{hurt}; in
a three-seed replication the 30-qubit Rydberg block achieved the best mean and
lowest variance ($0.955 \pm 0.102$ vs $1.111 \pm 0.174$ baseline), with the
variant ordering analog $>$ digital $>$ none $>$ MLP. These pilot effects are
directional at current power; the full-protocol verdict --- positive or
negative --- fine-tunes the official checkpoint (which we reproduce at 0.1934
test VRMSE) and is reported at the workshop.
\end{abstract}

\section{Introduction}
\label{sec:intro}

World models --- learned simulators that roll a latent state forward under
dynamics --- increasingly serve as the interface between machine learning and
physics, from PDE surrogacy to video generation
\citep{ohana2024thewell,mccabe2025walrus,gao2026lingbotworld}. Their central
unsolved failure is \emph{physical consistency over long horizons}: on The
Well benchmark, the best official baseline (CNextU-net) reaches one-step test
VRMSE 0.1956 on \texttt{turbulent\_radiative\_layer\_2D} but degrades to 1.01
by rollout window 13:30 --- worse than predicting the spatial mean --- and even
the 1.3B cross-domain foundation model Walrus exceeds this floor at long
horizon \citep{ohana2024thewell,mccabe2025walrus}. Dedicated video-physics
benchmarks report the same gap for generative world models
\citep{motamed2025physicsiq,li2025worldmodelbench}.

Meanwhile, two literatures make claims that, if true, bear directly on this
failure. First, \emph{analog quantum simulation}: neutral-atom platforms
natively evolve Rydberg Hamiltonians, and reservoir-computing experiments on
QuEra Aquila encode inputs directly into detunings, positions, and drives of
up to 108 qubits \citep{kornjaca2024aquila}; Hamiltonian input encoding is
provably expressive, offering infinite-order nonlinearity in the encoded
variables \citep{xia2026magicqrc}. Second, \emph{quantum modules for
foundation models}: unitary and quantum(-inspired) adapters now fine-tune
frozen LLMs and vision--language models at extreme parameter efficiency
\citep{koikeakino2025quantumpeft,raj2025quic,aizpurua2026cayley,zhao2026mqadapter}.
Neither literature, however, has produced a matched-control win on a public
physical-simulation benchmark: the closest attempt, a 5-qubit mixer in a 3D
FNO surrogate, beat a larger classical model but \emph{lost to its own
parameter-matched control} \citep{papierz2026hqlpfno}, and reservoir studies
show apparent quantum edges evaporating under equal-tuning classical baselines
\citep{pandey2026qrcbaselines}. No quantum-hybrid entry exists on The Well,
PDEBench, or LagrangeBench.

Our premise connects the two claims. Analog quantum hardware may be the
natural substrate for physical dynamics, but no practical application fits on
today's qubit counts whole --- and real simulations are rarely uniformly
quantum-hard: they mix classically easy, effectively discrete structure with a
harder analog core. We therefore treat the qubit budget as a scarce resource
to be \emph{allocated}: a reconfigurable, input-parameterized register is
steered toward the highest-complexity analog subprocess of the problem at
hand, while the frozen world model carries the rest. If the division of labor
is right, small fixed budgets become useful on current hardware, and fidelity
versus qubit count becomes a measurable scaling question rather than a wait
for million-qubit machines. This is the hypothesis our protocol is built to
confirm or kill.

We take both claims seriously by making them falsifiable in one setting.
Our contributions, ordered by novelty:
\begin{enumerate}
  \item \textbf{The first quantum adapter inside a frozen video world
  model.} A quantum bottleneck at the KV-cache chunk boundary of a frozen
  14B open-weights interactive world model
  \citep{gao2026lingbotworld,wan2025wan}, scored by official public judges
  \citep{li2025worldmodelbench,motamed2025physicsiq} --- to our knowledge
  the first quantum module in any video world model and the first quantum
  entries on either benchmark. (The similarly-named QSWM
  \citep{sun2026qswm} is quantum-\emph{inspired} --- complex-valued,
  density-matrix-like latents evaluated on cellular automata --- and
  contains no quantum module.)
  \item \textbf{Checkpoint surgery against a published number.} We insert
  the bottleneck at the neck of the \emph{official released} CNextU-net
  checkpoint from The Well and fine-tune, targeting its published 0.1956
  one-step test VRMSE \citep{ohana2024thewell} with paired per-window
  Wilcoxon significance over the 873 official test windows. Quantum
  U-Net bottlenecks exist from scratch
  \citep{ma2025qufex,liu2026hqfnet,ali2026phqfno}; no quantum work has
  touched The Well or any released scientific-ML checkpoint --- the
  methodology (official-checkpoint surgery + published-number target +
  paired statistics) and its result are the contribution, not the
  architecture.
  \item \textbf{An input-conditioned analog ansatz, as a composite.} The
  Rydberg generator is parameterized by the latent state --- data-dependent
  detunings, gradient-trained 2D atom geometry, global drive --- via exact
  statevector autodiff inside the classical network, and the trained object
  is a valid neutral-atom program (Pulser-exportable). Geometry-as-parameter
  is the natively realizable knob of this platform: optical tweezers
  assemble atoms into essentially arbitrary, reconfigurable 2D/3D
  arrangements \citep{barredo2018synthetic}, now demonstrated at the
  6{,}100-qubit scale \citep{manetsch2025tweezer}, so a trained register
  layout is a deployable program, not an abstraction. Each ingredient
  exists in isolation \citep{silver2024recon,chen2026resq,pulserdiff2025};
  the composite, embedded as a layer of a predictive world model, does not.
  Thirty simulated qubits are realized as three blockade-decoupled
  10-qubit registers --- an implementation detail with prior art in patched
  quantum generators \citep{huang2020patchgan}.
  \item \textbf{A pre-registered, kill-rule protocol.} Qubit-scaling
  claims already exist in looser forms --- IonQ reports monotone scaling
  against parameter-matched classical baselines on a pretrained model
  \citep{ionq2026scaling}, and analog reservoirs scale to 108 qubits
  \citep{kornjaca2024aquila} --- so we do not claim ``more qubits helps''
  as novel. We contribute its first \emph{pre-registered,
  tuning-and-parameter-matched, significance-tested} instantiation on a
  public PDE benchmark: matched MLP/orthogonal/tensor-network/product-state
  controls \citep{pandey2026qrcbaselines,bowles2024better}, explicit kill
  rules (if controls match, the negative result is the paper), and a
  dequantization suite (bond-dimension sweeps, product-state controls,
  finite-shot gradients), under the honest-framing standard of
  \citet{schuld2022advantage} and the classical-simulability caveat of
  \citet{cerezo2025barren}.
\end{enumerate}

\section{Problem definition}
\label{sec:problem}

\paragraph{Task.} Let $x_t \in \mathbb{R}^{C \times H \times W}$ be a physical
field snapshot (here: density, pressure, velocity on a $128{\times}384$ grid).
Given a history $(x_{t-3},\dots,x_t)$, predict $x_{t+1}$; long-horizon
behavior is evaluated by autoregressive rollout. The official metric is
variance-scaled RMSE,
$\mathrm{VRMSE}(\hat x, x) = \sqrt{\|\hat x - x\|_2^2 \,/\,
(\|x - \bar x\|_2^2 + \epsilon)}$,
with $\mathrm{VRMSE}=1$ equivalent to predicting the spatial mean $\bar x$
\citep{ohana2024thewell}.

\paragraph{Hypotheses.}
\textbf{H1} (analog simulation): a quantum bottleneck with input-conditioned
Rydberg dynamics improves rollout accuracy over parameter- and tuning-matched
classical controls. \textbf{H2} (foundation models): the same block, as an
adapter in a frozen video world model, improves physical-consistency scores
over matched LoRA/MLP adapters. Each hypothesis carries a kill rule: if the
matched control is statistically indistinguishable from or better than the
quantum variant across seeds, the prong is reported dead in this regime.

\paragraph{Primary goal.} One-step test VRMSE $< 0.1956$ on
\texttt{turbulent\_radiative\_layer\_2D}, significant under a paired
per-sample Wilcoxon test over the $\sim$870 official test windows against the
official CNextU-net checkpoint, across 3 training seeds (Holm-corrected).

\paragraph{On WaveLiT.} Among compact architectures evaluated on TRL2D,
WaveLiT \citep{sankaran2026wavelit} reports one-step VRMSE 0.1308 (30.5M) and
0.1346 (9M) using the benchmark's 4-step history and stating consistency with
the official evaluation. The aggregation statistic, test split, and
normalization are, however, undocumented and no code is released; the authors'
extended version switches to median VRMSE (0.1167 for the same 9M model), a
statistic not comparable to the benchmark's mean convention, and the 30.5M
model's training budget exceeds the benchmark's 12h single-GPU convention by
roughly $2$--$3\times$. We treat WaveLiT as a promising but not
protocol-verified comparison point and evaluate exclusively under the official
mean-VRMSE protocol on the official test split.

\section{Method (summary)}
\label{sec:method}

A classical trunk produces a latent field $z$; an attention pool yields
$h \in \mathbb{R}^d$; a linear map sets per-qubit, per-layer detunings
$\delta = \pi \tanh(W h)$; each 10-qubit register evolves under
$H_{\mathrm{Ryd}}/\hbar = \sum_i \tfrac{\Omega_\ell}{2} X_i - \sum_i
\delta_{\ell,i} n_i + \sum_{i<j} C_6 |r_i - r_j|^{-6} n_i n_j$
with trainable times $\tau_\ell$, global drives $\Omega_\ell$, and trainable
coordinates $r_i$ (geometry-as-parameter); $\langle Z_i \rangle$ readouts
re-enter the trunk through a zero-initialized gated residual. The quantum
core is $\sim$500--2{,}100 parameters per register; every baseline implements
the identical pool$\to$core$\to$residual interface at matched total parameters
($\pm 0.3\%$). Training uses a batched complex64 statevector propagator with
exact autograd gradients, cross-checked against an independent simulator to
$10^{-6}$. Prong H2 taps the denoised chunk latent of the frozen world model
at its KV-cache chunk boundary, one adapter call per chunk.

\section{Pilot evidence (time-boxed sprints; internal protocol)}
\label{sec:pilot}

Pilot sprints use a stratified subset (3 of 9 cooling regimes, $2\times$
downsampled, single-frame input, 5-minute wall-clock or fixed-step budgets,
$\sim$1.33M parameters, fixed seeds) --- internally comparable, \emph{not}
comparable to the published table; the full-protocol runs target
Section~\ref{sec:problem}'s goal directly.

\begin{table}[h]
\centering
\small
\begin{tabular}{lcccc}
\toprule
Variant & VRMSE@1 & VRMSE@8 & steps & note \\
\midrule
persistence floor & 0.583 & 1.002 & --- & trivial baseline \\
FNO trunk (none) & 0.419 & 1.143 & 10{,}597 & \\
$+$ matched MLP control & 0.425 & 1.213 & 10{,}326 & control \emph{hurts} \\
$+$ 12q HEA bottleneck & \textbf{0.404} & \textbf{0.876} & 6{,}912 & $-23\%$ rollout \\
$+$ 16q Rydberg-analog & 0.432 & 0.913 & 2{,}850 & beats floor \\
\midrule
fixed 2{,}500 steps: none & 0.477 & 2.102 & 2{,}500 & \\
fixed 2{,}500 steps: $+$12q & 0.486 & \textbf{1.325} & 2{,}500 & $-37\%$ per-step \\
\midrule
\multicolumn{5}{l}{\emph{3-seed study (mean$\pm$std over seeds):}} \\
none & --- & $1.111 \pm 0.174$ & $\sim$10.3k & \\
$+$ matched MLP (30-dim) & --- & $1.123 \pm 0.211$ & $\sim$10.1k & control \\
$+$ 30q HEA ($3{\times}10$q) & --- & $1.063 \pm 0.152$ & $\sim$4.4k & \\
$+$ 30q Rydberg ($3{\times}10$q) & --- & $\mathbf{0.955 \pm 0.102}$ & $\sim$5.4k & best mean \& lowest var \\
\bottomrule
\end{tabular}
\vspace{2pt}
\caption{Pilot sprints (The Well TRL2D subset). Only quantum variants beat
the persistence floor at 8-step rollout; the effect persists at fixed step
budget (throughput removed). Multi-seed replication is in progress: run-to-run
$\sigma(\mathrm{VRMSE@8}) \approx 0.2$ rivals single-run effect sizes, so
per-regime paired statistics (9 points/comparison) gate any claim.}
\label{tab:pilot}
\end{table}

Across three fixed-seed replicates, the variant ordering matches the
hypothesis pattern --- analog $>$ digital $>$ no bottleneck $>$ classical
control --- with the 30-simulated-qubit Rydberg block best on both mean and
variance, and the analog qubit-scaling trend monotone in our runs
(12q: 1.17 $\to$ 16q: 0.91 $\to$ $3{\times}10$q: $0.955 \pm 0.102$). With
$n{=}3$ aggregate seeds these differences are directional, not significant;
the per-regime paired protocol and the full-protocol runs below carry the
statistical burden. As the anchor for the primary goal, we re-evaluated the
official CNextU-net checkpoint on the official test split with our pipeline:
mean VRMSE \textbf{0.1934} over 873 windows (published: 0.1956; 1.1\%
discrepancy), and we adopt our stricter reproduced number as the bar.

\section{Results}
\label{sec:results}

Four result blocks, in decreasing order of maturity: a completed
three-seed internal study whose numbers are final
(\S\ref{sec:wave15}); the official-benchmark campaign, in flight at
submission-scaffolding time, whose cells are visibly marked placeholders
(\S\ref{sec:campaign}); released artifacts (\S\ref{sec:artifacts}); and the
full evidence ladder with honest per-benchmark status (\S\ref{sec:ladder}).

\subsection{Sprint arena, wave 1.5: rollout-aware training unlocks the
quantum edge}
\label{sec:wave15}

The decisive internal experiment crosses the training objective with the
bottleneck variant at a fixed 2{,}000-step budget (throughput removed), 3
seeds: one-step MSE ($k{=}1$) versus a $k{=}4$ pushforward/TD objective
(inner rollout steps stop-gradiented, loss on the final step). Results in
Table~\ref{tab:wave15} are complete and final.

\begin{table}[h]
\centering
\small
\begin{tabular}{lccc}
\toprule
Variant & VRMSE@8, $k{=}1$ & VRMSE@8, $k{=}4$ (TD) & TD gain \\
\midrule
FNO trunk (none) & $0.964 \pm 0.005$ & $0.957 \pm 0.255$ & $\approx 0\%$ \\
$+$ matched MLP (control) & $1.057 \pm 0.057$ & $0.953 \pm 0.135$ & $-10\%$ \\
$+$ 30q Rydberg ($3{\times}10$q) & $1.100 \pm 0.194$ & $\mathbf{0.845 \pm 0.051}$ & $\mathbf{-23\%}$ \\
$+$ 30q Rydberg, conditioned geometry & $1.043 \pm 0.210$ & $0.895 \pm 0.137$ & $-14\%$ \\
\bottomrule
\end{tabular}
\vspace{2pt}
\caption{Wave 1.5 (The Well TRL2D subset, internal protocol): 8-step rollout
VRMSE, mean$\pm$std over seeds \{1337, 1338, 1339\}, fixed 2{,}000 steps.
Under the TD objective the 30-simulated-qubit Rydberg bottleneck beats the
parameter-matched MLP control on \textbf{3/3 seeds} (0.903/0.826/0.806 vs
1.098/0.933/0.829) and posts the best mean and lowest variance overall.
Interaction effect: switching to the TD objective improves the quantum
variant by $-23\%$, the classical control by $-10\%$, and the bare trunk not
at all --- under one-step training the trunk alone is best and the quantum
edge is \emph{invisible}. Adapters need rollout-aware training; this finding
sets the round-2 design in \S\ref{sec:campaign}.}
\label{tab:wave15}
\end{table}

This is the pre-registered H1a pattern (quantum $<$ both controls) at its
strongest to date, and it explains the flat one-step fine-tunes observed in
early checkpoint-surgery pilots: the quantum bottleneck's contribution lives
in rollout dynamics, which the one-step loss never sees.

\subsection{Official-benchmark campaign (TRL2D): checkpoint surgery in two
rounds}
\label{sec:campaign}

The campaign fine-tunes the \emph{official released} CNextU-net checkpoint
under the official protocol, three arms (vanilla continued / matched-MLP
control / 30q Rydberg bottleneck) $\times$ three seeds, with single-shot
test evaluation per best-on-validation model and paired per-window Wilcoxon
against our reproduced anchor (\TRLanchorRepro{} over \TRLwindows{} windows;
published: \TRLanchorPub). Round 1 trains with the benchmark's one-step
objective; round 2 repeats the surgery with the wave-1.5-informed TD
objective ($k{=}4$). A 3k-step pilot preserved the hypothesis ordering on
validation VRMSE (quantum $0.18977 <$ vanilla $0.18985 <$ MLP $0.19009$),
differences within noise --- consistent with wave 1.5's finding that
one-step training barely moves any arm. Numbers below are placeholders
until the campaign's paired statistics land; the verdict is reported either
way. \CampaignVerdict

\begin{table}[h]
\centering
\small
\setlength{\tabcolsep}{5pt}
\begin{tabular}{llccc}
\toprule
Round & Arm & Test VRMSE (3 seeds) & Wilcoxon $p$ vs anchor & win rate \\
\midrule
--- & official checkpoint (anchor) & \TRLanchorRepro & --- & --- \\
\midrule
1 (one-step) & vanilla continued & \TRLvanilla & --- & --- \\
1 (one-step) & $+$ matched MLP (control) & \TRLmlp & --- & --- \\
1 (one-step) & $+$ 30q Rydberg & \TRLquantum & \TRLpvalue & \TRLwinrate \\
\midrule
2 (TD, $k{=}4$) & vanilla continued & \TDvanilla & --- & --- \\
2 (TD, $k{=}4$) & $+$ matched MLP (control) & \TDmlp & --- & --- \\
2 (TD, $k{=}4$) & $+$ 30q Rydberg & \TDquantum & \TDpvalue & \TDwinrate \\
\bottomrule
\end{tabular}
\vspace{2pt}
\caption{TRL2D official-protocol campaign (scaffolding; red cells are
in-flight placeholders). One-step test mean-VRMSE on the official test
split; paired one-sided Wilcoxon (arm $<$ anchor) over \TRLwindows{}
per-window scores, Holm-corrected across seeds. Pass rule (pre-registered
H1d): quantum below both controls \emph{and} below the \TRLanchorRepro{}
anchor at $p<0.05$. Kill rule: if the matched control is indistinguishable
from or better than the quantum arm, the prong is reported dead in this
regime and this table ships with that verdict.}
\label{tab:campaign}
\end{table}

\subsection{Released artifacts}
\label{sec:artifacts}

\paragraph{TORAX-Bench v1.} With this paper we release a new benchmark:
1{,}000 actuator-conditioned tokamak core-transport discharges generated
with the TORAX simulator (all four transport equations evolved --- ion and
electron heat, current, electron density), sampled i.i.d.\ from a
12-dimensional scenario space with \emph{time-varying} actuators
(plasma-current ramp, three-knot heating waveform, gas-puff step). Each
discharge saves $[T_e, T_i, n_e, q]$ profiles on a uniform 121-frame,
6-second grid plus the commanded actuator waveforms; splits are frozen at
800/100/100 with a dataset card and full generation manifest (44\,MB
total). Reference CPU baselines are set and frozen: persistence@8 mean
rollout VRMSE \textbf{0.0923} (the number to beat) and per-timestep
climatology 0.5546. Adapter numbers on this benchmark are pending and not
claimed here.

\paragraph{Pre-registered protocol and reproduction pack.} The second
artifact is the protocol itself: the hypothesis scoreboard with explicit
pass/kill rules per claim, the per-run results ledger, the paired-statistics
scripts, and a \texttt{REPRODUCE.md} that maps every number appearing in
this paper --- including every pending cell above --- to the exact command or
executed notebook that regenerates it. All runs execute on a single
126\,GB workstation.

\subsection{Evidence ladder}
\label{sec:ladder}

Table~\ref{tab:ladder} is the full claim surface with per-benchmark status
at scaffolding time --- filled, in-flight, and not-yet-started rows alike.
Rows convert from ``pending'' to numbers only through the pre-registered
commands; nothing is claimed for a pending row.

\begin{table}[h]
\centering
\scriptsize
\setlength{\tabcolsep}{3pt}
\begin{tabular}{lllll}
\toprule
Benchmark & Metric & Reference point & Ours & Status \\
\midrule
The Well TRL2D (official) & VRMSE $\downarrow$ & CNextU-net \TRLanchorPub{} / \TRLanchorRepro & Table~\ref{tab:campaign} & campaign in flight \\
TRL2D sprint arena (internal) & VRMSE@8 $\downarrow$ & best classical $0.953 \pm 0.135$ & $\mathbf{0.845 \pm 0.051}$, 3/3 seeds & \textbf{complete} \\
The Well MHD\_64 & VRMSE $\downarrow$ & CNextU-net 0.1633 & --- & pipeline ready \\
TORAX-Bench v1 (ours) & VRMSE@8 $\downarrow$ & persistence@8 0.0923 & baselines released & adapter pending \\
Physics-IQ & score $\uparrow$ & Cosmos $\approx 63.4$ & --- & pipeline staged \\
WorldModelBench & judge score $\uparrow$ & leaderboard & --- & queued (full version) \\
TokaMark G2 (MAST magnetics) & NRMSE $\downarrow$ & CNN+LSTM 0.1343 & --- & data access verified \\
Sophelio equilibrium challenge & $G_{\mathrm{ratio}}$ $\uparrow$ & Ch2 open $\approx 0.654$ & --- & committed (Oct 2026) \\
Pelzer 1HP heat plume & MAE $\downarrow$ & U-Net 0.046\,$^{\circ}$C & --- & loader + split frozen \\
\bottomrule
\end{tabular}
\vspace{2pt}
\caption{Evidence ladder: one adapter stack, every benchmark it is
pre-registered against, and the honest status of each row. ``Reference
point'' is the published or strongest available number for that benchmark;
dashes are not results. The two filled rows are the sprint-arena study
(Table~\ref{tab:wave15}) and the TORAX-Bench baseline release
(\S\ref{sec:artifacts}).}
\label{tab:ladder}
\end{table}

\section{Roadmap and limitations}
\label{sec:roadmap}

Full-protocol runs fine-tune the official CNextU-net checkpoint with and
without the quantum bottleneck under identical budgets. Known limitations we
carry openly: the pooled register is a poly($n$) information bottleneck
(readout scaling is reported, not hidden); entangled 30-qubit training uses
bounded-$\chi$ tensor-network simulation (exact statevector at 30q is
inference-only on our hardware); sprint-scale variance mandates the paired
protocol above; and all quantum components are classically simulated, so
hardware execution (via Pulser export to neutral-atom platforms) is future
work, not a claim. All experiments, including every negative result, run on a
single 126\,GB workstation and reproduce from one command per table.

\bibliographystyle{plainnat}
\bibliography{references}

\end{document}
